\section*{The Propers: Detailed Commentary}

\subsection*{Help that reforms the one who asks\enspace\properrefs{Int., Coll., All.}}

Psalm 53 (54) begins with threatened prayer, not retrospective calm. Its title places David within the Ziphite reports to Saul, but that inherited horizon is not a recovered composition scene. The appointed antiphon moves from plea to recognition: God helps and sustains while danger remains, and rescue and truthful judgment are still entrusted to him. Theodoret reads that turn as foreknowledge of help soon to come and the resulting praise as a willing sacrifice. Augustine reads divine help as interior cooperation rather than replacement of the servant's willing action (\emph{Exposition on Psalm 54}, 4--5, 8). Hilary emphasizes confidence under persecution, while Bellarmine makes the antiphon's ``behold'' a turn toward present assurance and willing praise. Their shared point is active dependence.

The Collect specifies what such dependence does to desire. Its grammar does not ask God first to satisfy the petitioners and afterward to approve their wants. It asks him to make them request what pleases him so that what they desire may rightly be granted. Augustine's \emph{Letter 130 to Proba}, \S\S~15--18 and 21--29, does not comment on the prayer but illuminates its doctrine: the Lord's Prayer gathers desire into Christ's form, delay enlarges capacity for the gift, and refusal of a harmful temporal request may itself be mercy. The Collect thus seeks conversion of the desiring subject, not a guarantee for every sincere request.

Psalm 58 (59) returns the same dependence before the Gospel. Its title recalls Saul's agents watching David's house (1~Sam. 19:11), so the peril is literal before it is metaphorical. Theodoret first calls the petition just because it seeks release from attackers who wage an unjust war. Bellarmine similarly begins with David and then widens the prayer to any just person in mortal danger; Augustine redirects it toward pride, spiritual enemies, and temptation. Taken together, the three readings are complementary: they hold the historical sufferer, legitimate protection, and divine judgment in one view.

\begin{threestable}{Appointed movement}{Interpretive force}{Precise sense}
\textbf{God helps} & Aid is already present within the servant's prayer and action. & Dependence neither erases agency nor guarantees immediate rescue.\\
\textbf{God forms desire} & Mercy orders the petitioner before granting the petition. & Intensity and sincerity do not make every object good.\\
\textbf{God delivers} & Real danger and spiritual temptation are entrusted to divine judgment. & Protection and judgment remain entrusted to God.\\
\end{threestable}

\sectionguard
\subsection*{Sacramental privilege under judgment\enspace\properrefs{Ep., Grad.}}

The lesson begins in the middle of Paul's answer about idol-associated meals. In 1~Corinthians 10:1--5 the sea and cloud, spiritual food and drink, and Christ the rock make Israel's history urgently relevant to baptized Christians. Verses 6--10 then move through craving (Num.~11), the idol feast (Exod.~32:6), sexual sin (Num.~25), testing and serpents (Num.~21), and murmuring with the destroyer. The missal prints Paul's twenty-three thousand at verse 8; Numbers 25:9 gives twenty-four thousand. Aquinas discusses the inherited difference, and the lesson keeps Paul's figure.

The argumentative center is addressed to ``us.'' The events truly occurred, Paul says, and were written for Christian correction; the person who thinks he stands is the one commanded to take heed. Chrysostom's \emph{Homily 23 on First Corinthians} treats the catalogue as a deliberate dismantling of Corinthian complacency. Aquinas, \emph{Super I Corinthios} X, lect.~2, distinguishes sins of heart, deed, and speech within Israel's real history. Romans 9 and 11 name Israel's gifts, covenants, promises, and calling while warning Gentile Christians against boasting.

Ambrosiaster locates verse 12 even more narrowly within the idol-meal dispute. Those presuming on their knowledge that all foods were lawful had become a scandal to weaker siblings; supposed progress under rival teachers had actually made them worse. The warning therefore tests neither private confidence nor doctrinal knowledge in isolation, but their effect on another member. At verses 16--17 Ambrosiaster makes the consequence explicit: sharing one bread and cup means that members of one body should share one mind and action.

Verse 13 consoles without reversing the warning. Chrysostom reads temptation as genuinely human yet endurance as impossible without God's aid; Aquinas analyzes temptation, consent, ready help, grace, and constancy rather than autonomous capacity. The promise concerns temptation to sin and addresses a community in the plural. It is not the slogan that every trauma will remain within one isolated person's strength, nor does a need for help reveal defective faith.

Paul's next word is ``therefore'': flee idolatry (10:14). He then interprets the Eucharistic cup and bread as participation in Christ and concludes that the many are one body because they share one bread (10:16--17). Chrysostom presses both union with Christ and the ethical contradiction of division; Aquinas, lect.~4, similarly moves from sacramental participation to ecclesial unity. That canonical continuation creates a strong bridge to this Sunday's Communion and Postcommunion.

Psalm 8 places the warned assembly beneath the Lord's admirable name and majesty. Its full context joins creaturely smallness to received dignity, and Hebrews 2 later reads its human vocation christologically. In the Greek tradition Theodoret takes ``how admirable'' intensively: God's name is acclaimed everywhere as that of heaven and earth's creator, while his providence reaches fragile humanity. Augustine develops the christological ascent; Aquinas emphasizes divine excellence manifested throughout earth and above heaven; Bellarmine moves from created beauty to a magnificence creation cannot contain. The Gradual therefore answers presumed standing with wonder, not self-contempt.

\sectionguard
\subsection*{The visitation recognized through tears\enspace\properrefs{Gosp.}}

Luke's sequence controls the interpretation. The royal approach has just acclaimed the king and ``peace in heaven'' (19:38). Jesus then sees Jerusalem, weeps, laments the things of peace now hidden, and foretells enemies enclosing the city because its visitation was not recognized. He next enters the Temple, corrects it, and remains there teaching daily. The tears, siege oracle, Temple action, and continued teaching form one narrated movement of compassion, judgment, reform, and instruction.

Josephus supplies a historical comparison. \emph{Jewish War} V.12.1--4 describes Titus's circumvallation around Jerusalem, the cutting off of escape, famine, and whole households dying; VII.1.1 describes the city's and Temple's demolition while naming spared towers and part of the western wall. The correspondences make Luke's siege language concrete and preserve the victims' human catastrophe. They do not settle the Gospel's composition date, and they do not show that literally every masonry stone disappeared.

The two prophetic citations deepen the Temple action. Isaiah 56:1--8 joins justice and covenant fidelity to the inclusion of foreigners and eunuchs, acceptable sacrifice, a house of prayer for all peoples, and God's gathering of others to those already gathered. Jeremiah 7:1--15 attacks false confidence in the sanctuary while worshipers steal, murder, oppress the alien, orphan, and widow, shed innocent blood, and pursue idols. A ``den of robbers'' is therefore not merely a market inside sacred architecture; it is a refuge claimed by unreformed injustice. Jesus' action is intra-Jewish prophetic correction ordered to inclusive prayer and truthful life.

The Gospel's direct expositors make genuinely different moves. Origen's \emph{Homily 38 on Luke}, in Jerome's Latin translation, sees Jesus ground the beatitude of mourning by his own tears. He then asks whether the lament reaches Christians after Gospel teaching and sacramental initiation; his Christian Jerusalem can still fail to recognize visitation. The homily next condemns selling Spirit-given teaching and ends at the Savior's banquet, the bread of life, and Christ's flesh and blood. It therefore supplies an ancient Gospel-to-Communion movement, made within his exposition of Luke. Cyril of Alexandria independently stresses that Christ's tears manifest pity and a will for peace and salvation.

Ambrose reads the Temple not as a merchant's lodging but as holiness's dwelling, priestly ministry as gratuitous service, Scripture as divine coin requiring discernment, and the sale of doves as traffic in the Spirit's grace. Gregory the Great turns visitation toward hearers lulled by prosperity, then reads sacred office and the conscience as Temples that must become places of prayer and daily teaching. Bonaventure presents the tears as priestly compassion even on a day of royal honor, connects the Temple citations to avarice and simony, and emphasizes public teaching after correction. Anthony of Padua spiritualizes Jerusalem as the soul whose ruined virtues call for contrition. These applications turn the Gospel toward Christian self-examination while keeping the historical city and its suffering in view.

After expounding Luke 19:41--48 verse by verse, Bede adds a Paschal synthesis. He aligns the five days from Christ's reception near the Mount of Olives to Passover with Exodus 12's keeping of the victim, then moves from the disciples' attentive reception through the gift of Christ's Body and Blood to his arrest and death. This direct Gospel-to-Passion-and-Eucharist bridge is Bede's reception of Luke's sequence.

The Pontifical Biblical Commission, \emph{The Jewish People and Their Sacred Scriptures in the Christian Bible}, no.~74, reads Luke's tears together with 23:28--31: Jesus' attention rests on the coming suffering of Jerusalem's women and children even as his own Passion approaches. Jesus, his disciples, Jerusalem, and the Temple all stand within Israel. Within this horizon, the Christian assembly receives the warning as a summons to recognize peace, protect the vulnerable, reform worship, and remain teachable.

\sectionguard
\subsection*{The word that delights, warns, and is kept\enspace\properrefs{Off.}}

Psalm 18 (19) has already moved from the heavens' wordless witness to the Lord's instruction before the Offertory begins. The appointed span keeps together right precepts that gladden, a command that enlightens, filial fear that endures, true judgments, value beyond gold, sweetness beyond honey, warning, and the servant's keeping. The selected antiphon is therefore neither pleasant mood nor praise of judgment abstracted from obedience.

Theodoret keeps the Greek terms distinct: precepts gladden by showing the cause of justification, the command illumines the mind's sight, and judgments declare divine decisions; their sweetness belongs to a life ordered as truly human. Augustine's \emph{Exposition on Psalm 19}, 9--12, interprets right judgments as ordering the interior person and their sweetness as delight beyond bodily taste. Aquinas's \emph{Super Psalmo 18} develops the sequence as right action, illumination, enduring chaste fear, truth, warning, and observance. Bellarmine adds the experiential claim that the servant discovers fruit in keeping the judgments. The witnesses agree that delight is an appetite schooled by truth; they differ in emphasis among interior order, moral taxonomy, and experienced practice.

Heard after Luke, the Offertory prevents two reductions. Judgment cannot mean destruction alone, because the servant delights in and keeps it. Sweetness cannot mean reassurance alone, because the same word warns and searches. The juxtaposition supports reform of worship, religious commerce, and treatment of the vulnerable, though the psalm does not predict the Temple scene.

\sectionguard
\subsection*{From the memorial to abiding in one body\enspace\properrefs{Sec., Comm., Postcomm.}}

The Secret's syntax is causal. The Church asks worthily to frequent the mysteries \emph{because}, whenever the memorial of this sacrifice is celebrated, the work of redemption is enacted. Its strongest reception quotes it directly. In \emph{Summa theologiae} III, q.~83, a.~2, objection 5 quotes this Sunday's opening clause; Thomas's answer rejects the objection's inference that frequency means many celebrations by one priest in a day and instead explains daily celebration by the daily need for the once-for-all Passion's fruit and perpetual memorial. The event in Christ the head occurred once; its fruit reaches the faithful sacramentally.

Jordi Pinell's textual history distinguishes the appointed 1962 wording from older witnesses. The Veronese has \latin{frequentata mysteria} and \latin{exeritur}; the Old Gelasian reshapes the petition toward \latin{frequentare}, the mixed Gregorian tradition receives it, and \latin{digne} appears in the mid-eleventh-century Rossianum. Pinell argues that the older verb presents redemption's work becoming manifest and operative in Christian life, probably within a Paschal setting. That historical proposal illuminates the prayer's depth, and the 1962 Missal still reads \latin{exercetur}.

The Second Vatican Council then makes the prayer programmatic. \emph{Sacrosanctum Concilium} 2 quotes \latin{opus nostrae Redemptionis exercetur}, and its first footnote explicitly identifies the Secret of the Ninth Sunday after Pentecost. \emph{Presbyterorum ordinis} 13 reuses the clause as \latin{opus nostrae redemptionis continuo exercetur} for Eucharistic sacrifice and again names this Sunday in its note; the added \latin{continuo} means that this is not a verbatim repetition. Most decisively for the complete formulary, \emph{Lumen gentium} 3 says that whenever the sacrifice of the Cross is celebrated on the altar, redemption's work is enacted; its next sentence says the Eucharistic bread represents and effects the unity of the faithful who are one body, citing 1~Corinthians 10:17. This is clear verbal reuse of the prayer's clause and a magisterial joining of sacrifice to unity. Unlike the other two acts, however, \emph{Lumen gentium} does not identify this Sunday in a note, so its link to the Secret rests on the wording alone.

\Needspace{10\baselineskip}
\begin{threestable}{Witness}{Interpretive move}{Distinction}
\textbf{Thomas} & Daily need receives the fruit of the Passion through perpetual memorial. & His quoted objection does not authorize repeated daily celebration by one priest.\\
\textbf{\emph{Sacrosanctum Concilium} / \emph{Presbyterorum ordinis}} & The Sunday Secret names what Eucharistic liturgy and priestly ministry enact. & The Church does not manufacture or repeat the historical Passion.\\
\textbf{\emph{Lumen gentium}} & Altar celebration of the Cross's sacrifice leads immediately to effected one-body unity. & Unity is Christ's sacramental gift and an ecclesial demand, not mere sentiment.\\
\end{threestable}

The Communion supplies the mode of participation: eating Christ's flesh and drinking his blood are joined to mutual abiding. Augustine's \emph{Tractate 26 on John} first rereads 1~Corinthians 10:1--4 to distinguish visible sacramental reception from fruitful participation, then interprets the one bread and one body as a sign of unity and bond of charity; his grains and grapes make incorporation communal. Chrysostom's \emph{Homily 47 on John} emphasizes Christ's self-gift and the life drawn from him. Cyril grounds that life in Christ's life-giving humanity united to the living Word. Aquinas, \emph{Super Ioannem} VI, lect.~7, makes mutual abiding the union and life to which reception is ordered. These complementary moves exclude a momentary bodily transaction without denying sacramental objectivity.

The Postcommunion states the test with unusual economy: communion in the sacrament is asked to confer purification and bestow unity. Personal cleansing that entrenches division contradicts the prayer's direction; institutional unity without conversion empties its first petition. The printed Trinity Preface keeps the action doxological, while the Pauline continuation keeps it one-bodied. The formulary thus ends not with the confidence of one who thinks he stands, but with received abiding whose fruit must become both a purified person and a united Church.

Dom Lucien Fromage, continuing Guéranger's \emph{The Liturgical Year}, comments directly on this complete liturgical close. He calls the Communion the mystery of divine union and identifies in the Postcommunion both the sanctification of each member and the unity of the social body (\emph{The Liturgical Year}, Time after Pentecost II, pp.~255--256).

Later Roman reception makes the Postcommunion's two petitions mutually interpretive. A 1970 Prayer over the Offerings for Christian unity borrows \latin{purificationem conferat} from this ancient prayer. Innocent Smith traces the reuse alongside \emph{Lumen gentium} 15, where the Church's purification and renewal serve a clearer sign of Christ and restored unity. The reuse is later than the 1962 Missal. Chant reception gives the Communion the same relational emphasis: Johner notes that a falling fourth on \latin{in me manet} is answered by a rising fourth on \latin{ego in eo}.

\begin{fourstable}{Proper}{Divine action}{Requested fruit}{What it is not}
\textbf{Secret} & Redemption's work is enacted in sacramental memorial. & Worthy, repeated participation. & No second historical crucifixion.\\
\textbf{Communion} & Christ gives himself and mutually abides. & Life and incorporation. & No merely momentary or private possession.\\
\textbf{Postcommunion} & Sacramental communion acts in the recipients. & Purification and unity. & Neither emotion nor sociology exhausts the fruit.\\
\textbf{Trinity Preface} & The appointed preface directs praise to the triune God. & Doxological worship. & It is liturgical context, not an eleventh proper.\\
\end{fourstable}
